
% EXAMPLES OF NEW COMMANDS
\newcommand{\bea}{\begin{eqnarray}} % Shortcut for equation arrays
\newcommand{\eea}{\end{eqnarray}}
\newcommand{\e}[1]{\times 10^{#1}}  % Powers of 10 notation

%----------------------------------------------------------------------------
%	ADD YOUR DEFINITIONS AND COMMANDS (be careful of existing commands)
%----------------------------------------------------------------------------

%->New theorem environments
%\newtheorem{problem}{Problem}
%\newtheorem{lemma}{Lemma}
%->New Math operators
%argmax, argmin
\DeclareMathOperator*{\argmax}{arg\,max}
\DeclareMathOperator*{\argmin}{arg\,min}
\newcommand{\vect}[1]{\boldsymbol{\mathbf{#1}}}
\newcommand{\norm}[1]{\| #1 \|}
\DeclareMathOperator{\Tr}{Tr}
\newcommand{\trace}[1]{\Tr(#1)}
\newcommand{\outbrackets}[1]{\langle #1 \rangle}
\newcommand{\outcurlybrackets}[1]{\{#1\}}
\DeclareMathOperator{\rk}{rank}
\newcommand{\rank}[1]{\rk(#1)}
\DeclareMathOperator{\vectorization}{vec}
\newcommand{\outbracket}[1]{\langle #1 \rangle}

\newcommand{\cind}{\mathrel{\perp\mspace{-9mu}\perp}}

\newtheorem{approach}{Approach}

% citations
\newcommand{\iconcite}[1]{%
  \raisebox{-0.3ex}{\pgfuseimage{beamericonarticle}}~\citealp{#1}%
}